\section{Introduction}
\label{sec:introduction}

A trusted research environment (TRE) lets researchers analyze sensitive data while controlling what they can take out~\cite{fivesafes,hdruk2020tre}. As local language models begin to do analytical work inside TREs, a TRE has to decide which models to admit. A companion system contains the analyst: it sees only aggregates that have passed disclosure control, and it has no network access~\cite{wade2026safetre}. A companion evaluation measured how often candidate analysts draw a wrong conclusion when the data invite one, and recommended a configuration, Qwen3.5 27B followed by a gpt-oss 20B checker, that met its release threshold on two sealed task sets~\cite{wade2026opacity}.

Those measurements assume an honest but fallible analyst, and a harder question follows for an open-weight model of any origin: how do we know that it does not mislead on purpose? No test can show that a model is trustworthy in general, because a trigger that was never tried would pass. Tests can characterize the behavior a model shows under specified conditions. They belong beside containment, pinned provenance and monitoring in service, which hold even when a test misses something.

This paper asks two questions of a candidate model before admission. The first is whether its internal state reveals errors that its own verdict misses: classifiers trained on a model's internal activations, called probes, could check a write-up from the inside if the state holds knowledge that the verdict leaves out. The second is whether a biased write-up can be detected. We plant a bias with a hidden instruction to steer the conclusion without changing any number, test whether probes, text classifiers and correctness checkers can tell a biased write-up from its honest pair on unseen tasks, and apply the resulting detector to the recommended analyst.

The paper makes three contributions. It shows that a probe's apparent gain over a model's verdict can come entirely from knowing which tasks are risky, and that a measure within tasks, with a probe trained on shuffled labels as a control, exposes this. It builds a positive control for instructed bias, with a check that the instruction worked, and finds that the text carries the bias as well as the internal state does, while checkers of correctness miss most of it. It turns that control into a bias test for a candidate model, shows that an ordinary change in the style a researcher asks for can break a detector's threshold, and applies the corrected test to the recommended analyst. The stop rules of all four experiments were fixed before their data were read. None of the experiments involves a model trained to deceive, although an engineering pilot prepares for that experiment.
